\section{Metodología}
\label{sec:methodology}

El estudio adopta un enfoque cuantitativo, analítico-experimental y computacional.
Su lógica es hipotético-deductiva: se formulan modelos y propiedades bajo
supuestos explícitos, se derivan predicciones contrastables y se evalúan mediante
demostraciones, enumeraciones exhaustivas y experimentos pareados de simulación.
El análisis formal determina qué puede garantizarse dentro de cada modelo; la
experimentación cuantifica rendimiento, robustez y fallos cuando la complejidad
impide obtener una caracterización cerrada.

La investigación no trata la coordinación multi-robot como un único bloque. Una
misión de transporte cooperativo enlaza decisiones de asignación, cierre de
coaliciones, compatibilidad de recursos, factibilidad mecánica, movimiento,
seguridad, reparación y comunicación. Evaluar toda la cadena desde el inicio
impediría atribuir un fallo a una causa concreta. Por ello, el problema se descompone
en los subproblemas SP0--SP8. Cada escalón conserva lo ya validado y añade una
condición nueva. La misión Cargo integra después las capas principales y el piloto
Industrial 2 comprueba su transferencia geométrica y cinemática a un escenario
más congestionado.

La modalidad física principal es Cargo planar soportado. Tras el acoplamiento, los
robots y la carga se representan mediante un cuerpo compuesto con contactos
fijos y esfuerzos acotados. El empuje unilateral y el \emph{caging} requieren
modelos distintos de contacto, fricción y confinamiento, por lo que quedan fuera
del diseño confirmatorio.

\subsection{Diseño de investigación y lógica incremental}
\label{subsec:research_design}

El diseño combina dos tipos de evidencia. La evidencia formal comprende
demostraciones de potencial, caracterizaciones de equilibrio, condiciones de
factibilidad, cotas y propiedades de terminación. Sus conclusiones son deductivas,
pero solo dentro de los supuestos enunciados. La evidencia experimental procede
de comparaciones controladas entre métodos que reciben el mismo escenario.
Permite estimar efectos y observar fallos, pero no sustituye una demostración de
convergencia, estabilidad o seguridad.

La unidad experimental independiente se denomina \emph{mundo}. Un mundo
\(w\) queda determinado por un escenario \(s_w\), un vector de factores
\(\zeta_w\) y una semilla pseudoaleatoria \(\sigma_w\):
%
\begin{equation}
    w = \left(s_w,\zeta_w,\sigma_w\right).
    \label{eq:world_definition}
\end{equation}
%
Todos los métodos de una comparación se ejecutan sobre el mismo mundo. Los
robots, cargas, contactos, mensajes e instantes de simulación están anidados
dentro de él y no se consideran réplicas independientes. Esta definición evita
aumentar artificialmente el tamaño muestral y conserva el emparejamiento entre
tratamientos.

La escalera incremental cumple cuatro funciones. Primero, separa propiedades que
no son equivalentes: una asignación factible puede ser mecánicamente irrealizable,
y una trayectoria segura puede bloquearse. Segundo, permite atribuir el efecto de
cada mecanismo mediante una ablación compatible. Tercero, mantiene oráculos y
comparadores dentro de tamaños en los que su solución puede certificarse.
Cuarto, impide transferir una garantía obtenida con poses fijas, contactos conocidos
o información global a una planta móvil con comunicación imperfecta.

\begin{figure}[htbp]
    \centering
    \begin{tikzpicture}[
        stage/.style={
            draw,
            rounded corners,
            align=center,
            minimum height=1.55cm,
            text width=0.205\textwidth,
            inner sep=4pt
        },
        integration/.style={
            draw,
            rounded corners,
            align=center,
            minimum height=1.35cm,
            text width=0.43\textwidth,
            inner sep=5pt
        },
        arrow/.style={-{Latex[length=2.5mm]}, thick}
    ]
        \node[stage] (decision) {
            \textbf{Decisión lógica}\\
            SP0--SP2\\
            asignación, cuotas\\
            y heterogeneidad
        };

        \node[stage, right=0.65cm of decision] (physics) {
            \textbf{Realización física}\\
            SP3--SP4\\
            wrench, acoplamiento\\
            y transporte
        };

        \node[stage, right=0.65cm of physics] (operation) {
            \textbf{Operación robusta}\\
            SP5--SP7\\
            obstáculos, fallos\\
            y tráfico
        };

        \node[stage, right=0.65cm of operation] (network) {
            \textbf{Información y escala}\\
            SP8\\
            radio, retardo, pérdida\\
            y tamaño de flota
        };

        \draw[arrow] (decision) -- (physics);
        \draw[arrow] (physics) -- (operation);
        \draw[arrow] (operation) -- (network);

        \node[integration, below=0.85cm of physics, xshift=2.25cm] (cargo) {
            \textbf{Integración Cargo}\\
            reclutamiento, acoplamiento, transporte,
            seguridad, fallo y reemplazo
        };

        \node[integration, below=0.55cm of cargo] (pilot) {
            \textbf{Piloto Industrial 2}\\
            prueba exploratoria de geometría, tráfico y bloqueo
        };

        \draw[arrow] (decision.south) |- (cargo.west);
        \draw[arrow] (physics.south) -- (cargo.north west);
        \draw[arrow] (operation.south) -- (cargo.north east);
        \draw[arrow] (network.south) |- (cargo.east);
        \draw[arrow] (cargo) -- (pilot);
    \end{tikzpicture}
    \caption{Lógica incremental de la validación. Cada bloque añade una
    condición que no puede inferirse del bloque anterior. La misión Cargo
    comprueba compatibilidad funcional, no una garantía híbrida global.}
    \label{fig:methodological_staircase}
\end{figure}

La Tabla~\ref{tab:sp_methodological_ladder} precisa la pregunta aislada en cada
escalón. La columna ``razón de separación'' identifica el error metodológico que
se produciría al omitirlo.

\begin{table}[htbp]
    \centering
    \caption{Escalera metodológica SP0--SP8 e integración final.}
    \label{tab:sp_methodological_ladder}
    \renewcommand{\arraystretch}{1.12}
    \begin{tabularx}{\textwidth}{
        >{\raggedright\arraybackslash}p{0.08\textwidth}
        >{\raggedright\arraybackslash}p{0.21\textwidth}
        >{\raggedright\arraybackslash}X
        >{\raggedright\arraybackslash}p{0.25\textwidth}
    }
        \toprule
        \textbf{Bloque} &
        \textbf{Condición añadida} &
        \textbf{Pregunta aislada} &
        \textbf{Razón de separación} \\
        \midrule

        SP0 &
        Asignación homogénea uno a uno &
        ¿El mecanismo alcanza cobertura y qué eficiencia pierde frente al
        problema de asignación exacta? &
        Fija oráculos, métricas y criterios de factibilidad antes de introducir
        coaliciones. \\

        SP1 &
        Cuotas y cardinalidad variable &
        ¿Cómo se forman grupos enteros y qué ocurre cuando la demanda supera
        la flota disponible? &
        Una preferencia continua no constituye una coalición ejecutable y el
        déficit lineal no selecciona grupos completos bajo escasez. \\

        SP2 &
        Capacidades y servicio heterogéneos &
        ¿Cómo cambia la selección cuando cada robot aporta una cantidad distinta
        a cada carga? &
        Cumplir una cardinalidad no garantiza que se cubra el servicio requerido. \\

        SP3 &
        Geometría, contactos y wrench &
        ¿La coalición seleccionada puede producir la fuerza y el momento
        exigidos? &
        La suma escalar de capacidades no representa direcciones de fuerza,
        brazos de palanca ni saturación. \\

        SP4 &
        Acoplamiento y dinámica de pose &
        ¿Los robots alcanzan los contactos y la carga sigue una referencia con
        entradas limitadas? &
        Un certificado estático no prueba aproximación, seguimiento ni
        estabilidad de la planta. \\

        SP5 &
        Obstáculos y huella compuesta &
        ¿La acción aplicada conserva separación sin impedir sistemáticamente el
        progreso? &
        El control de pose no garantiza seguridad y una guardia segura puede
        producir bloqueo. \\

        SP6 &
        Fallo simple y reserva &
        ¿Puede restaurarse el certificado y completarse el reemplazo antes del
        plazo disponible? &
        Reasignar recursos no implica que el reemplazo llegue ni que el contacto
        se restablezca a tiempo. \\

        SP7 &
        Tráfico entre coaliciones &
        ¿Cómo se excluyen conflictos y bloqueos entre cuerpos compuestos? &
        La seguridad frente a obstáculos estáticos no resuelve cruces, pasillos
        opuestos ni reserva de recursos compartidos. \\

        SP8 &
        Escala y comunicación imperfecta &
        ¿Qué calidad y coste conserva la coordinación con radio, retardo y
        pérdidas? &
        Un equilibrio calculado con información completa no representa lo que
        puede observar una red fragmentada u obsoleta. \\

        Cargo &
        Composición de las capas &
        ¿Pueden coexistir reclutamiento, acoplamiento, transporte, seguridad,
        fallo y reparación en una misión? &
        Verifica compatibilidad funcional, pero no convierte las garantías
        parciales en una garantía global. \\

        Industrial 2 &
        Geometría y tráfico más complejos &
        ¿Qué bloqueos aparecen fuera de los escenarios sintéticos simples? &
        Actúa como prueba exploratoria de transferencia, no como validación
        dinámica o industrial. \\

        \bottomrule
    \end{tabularx}
\end{table}

\subsection{Modelos y escalas de abstracción}
\label{subsec:models_abstraction}

Sean \(\mathcal{R}=\{1,\ldots,N\}\) el conjunto de robots y
\(\mathcal{L}=\{1,\ldots,K\}\) el conjunto de cargas. No todos los SP utilizan la
misma planta. SP0--SP2 fijan las poses para estudiar únicamente la decisión;
SP3 evalúa una carga estática y sus contactos; SP4--SP6 integran movimiento
planar; SP7--SP8 sustituyen la dinámica continua por rutas y recursos discretos.
La comparación dentro de cada SP conserva la misma escala de abstracción.

\paragraph{Robot diferencial.}
Durante la aproximación, cada robot se representa mediante el modelo de uniciclo.
Su pose es \(q_i=(p_i^x,p_i^y,\theta_i)\) y su entrada cinemática es
\(u_i=(v_i,\omega_i)\):
%
\begin{equation}
    \dot q_i =
    \begin{bmatrix}
        \cos\theta_i & 0 \\
        \sin\theta_i & 0 \\
        0            & 1
    \end{bmatrix}
    \begin{bmatrix}
        v_i\\
        \omega_i
    \end{bmatrix},
    \qquad
    |v_i|\leq v_i^{\max},
    \quad
    |\omega_i|\leq \omega_i^{\max}.
    \label{eq:unicycle_methodology}
\end{equation}
%
Para ruedas de radio \(r_w\) y semivía \(\ell_w\), las velocidades del uniciclo
se relacionan con las velocidades angulares izquierda y derecha mediante
%
\begin{equation}
    \begin{bmatrix}
        v_i\\
        \omega_i
    \end{bmatrix}
    =
    \frac{r_w}{2}
    \begin{bmatrix}
        1 & 1\\
        -1/\ell_w & 1/\ell_w
    \end{bmatrix}
    \begin{bmatrix}
        \dot\varphi_{i,L}\\
        \dot\varphi_{i,R}
    \end{bmatrix}.
    \label{eq:differential_drive_methodology}
\end{equation}
%
SP4 y Cargo incorporan dinámica de velocidad. Con
\(x_i=(q_i,v_i,\omega_i)\) y aceleraciones
\(a_i=(a_i^v,a_i^\omega)\),
%
\begin{equation}
    \dot x_i =
    \begin{bmatrix}
        v_i\cos\theta_i\\
        v_i\sin\theta_i\\
        \omega_i\\
        a_i^v\\
        a_i^\omega
    \end{bmatrix},
    \qquad
    a_i\in\mathcal{A}_i,
    \label{eq:dynamic_unicycle_methodology}
\end{equation}
%
donde \(\mathcal{A}_i\) recoge las cotas de aceleración y las limitaciones
derivadas de los actuadores. El modelo supone rodadura ideal y velocidad lateral
nula. No reproduce deslizamiento, deformación del neumático ni contacto
tridimensional rueda--suelo.

\paragraph{Asignación y cierre de coaliciones.}
La variable binaria \(x_{ik}\) indica si el robot \(i\) pertenece a la coalición de
la carga \(k\). La preferencia continua empleada por algunas dinámicas se denota
\(p_{ik}\in[0,1]\) y no se identifica con \(x_{ik}\). La coalición cerrada es
\(\mathcal{C}_k=\{i\in\mathcal{R}:x_{ik}=1\}\). Las restricciones lógicas comunes
son
%
\begin{equation}
    \sum_{k\in\mathcal{L}}x_{ik}\leq 1,
    \qquad
    |\mathcal{C}_k|\geq n_k,
    \qquad
    \sum_{i\in\mathcal{R}}c_i x_{ik}\succeq r_k,
    \qquad
    x_{ik}\in\{0,1\}.
    \label{eq:coalition_closure_methodology}
\end{equation}
%
La primera desigualdad impone exclusividad, \(n_k\) es la cardinalidad mínima y
\(c_i\succeq 0\), \(r_k\succeq 0\) son vectores de capacidad y requisito.
SP0 utiliza demanda unitaria; SP1 introduce \(n_k>1\); SP2 reemplaza la
equivalencia entre robots por contribuciones dependientes del par robot--carga.

\paragraph{Servicio operacional heterogéneo.}
La contribución usada en SP2 combina carga útil nominal, disponibilidad de
batería y distancia inicial. Para una reserva \(b^{\mathrm{res}}\), una escala de
distancia \(\ell_d\) y una referencia \(c^{\mathrm{ref}}>0\), se define
%
\begin{align}
    \psi_i^{b}
    &=
    \max\left\{
        0,
        \frac{b_i-b^{\mathrm{res}}}{1-b^{\mathrm{res}}}
    \right\},
    &
    \psi_{ik}^{d}
    &=
    \exp\left(-\frac{d_{ik}}{\ell_d}\right),
    \label{eq:availability_factors_methodology}\\
    a_{ik}
    &=
    \frac{c_i^{\mathrm{pay}}}{c^{\mathrm{ref}}}
    \psi_i^{b}\psi_{ik}^{d},
    &
    S_k(x)
    &=
    \sum_{i\in\mathcal{R}}a_{ik}x_{ik}.
    \label{eq:operational_service_methodology}
\end{align}
%
La carga \(k\) cubre su umbral operacional si
\(S_k(x)\geq d_k^{\mathrm{srv}}\). Esta condición mide servicio disponible; no
certifica fuerza, torque, soporte ni energía suficiente para toda la misión.

\paragraph{Factibilidad de wrench.}
Para una coalición y un conjunto de contactos fijos, \(\lambda_k\) agrupa los
esfuerzos admisibles y \(G_k(\mathcal{C}_k)\) es la matriz de agarre. El wrench
planar realizado es
%
\begin{equation}
    W_k = G_k(\mathcal{C}_k)\lambda_k,
    \qquad
    \lambda_k\in\mathcal{U}_k(\mathcal{C}_k).
    \label{eq:wrench_map_methodology}
\end{equation}
%
La factibilidad se evalúa mediante el residual normalizado
%
\begin{equation}
    \rho_k^{W}(\mathcal{C}_k)
    =
    \min_{\lambda_k\in\mathcal{U}_k(\mathcal{C}_k)}
    \left\|
        Q_W
        \left(
            G_k(\mathcal{C}_k)\lambda_k
            -W_k^{\mathrm{req}}
        \right)
    \right\|_2,
    \label{eq:wrench_residual_methodology}
\end{equation}
%
donde \(Q_W\) hace comparables fuerza y momento. Una coalición se acepta si
\(\rho_k^{W}\leq\varepsilon_W\). El problema es planar y cuasiestático; una
aceptación no demuestra estabilidad durante el transporte.

\paragraph{Carga planar compuesta.}
Después del acoplamiento, la pose de la carga es
\(q_k^L=(p_k^{L,x},p_k^{L,y},\theta_k^L)\) y su velocidad generalizada es
\(\nu_k=(v_k^x,v_k^y,\omega_k)\). La planta reducida satisface
%
\begin{equation}
    M_k\dot\nu_k+D_k\nu_k
    =
    G_k(\mathcal{C}_k)\lambda_k+W_k^{\mathrm{ext}},
    \qquad
    \dot q_k^L=T(\theta_k^L)\nu_k,
    \label{eq:payload_dynamics_methodology}
\end{equation}
%
con \(M_k\succ0\), \(D_k\succeq0\) y contactos fijos respecto del cuerpo. Los
esfuerzos se limitan antes de integrar. La reconstrucción rígida de los robots
alrededor de la carga no modela agarre, flexibilidad estructural ni pérdida de
contacto.

\paragraph{Seguridad geométrica.}
Sea \(\mathcal{F}_k(q_k^L)\) la huella compuesta de carga y robots y
\(\mathcal{O}_o\) un obstáculo. La holgura mínima se expresa como
%
\begin{equation}
    h_{ko}(q_k^L)
    =
    \operatorname{dist}
    \left(
        \mathcal{F}_k(q_k^L),\mathcal{O}_o
    \right)
    -d_{\mathrm{safe}}.
    \label{eq:safety_margin_methodology}
\end{equation}
%
El conjunto seguro nominal satisface \(h_{ko}\geq0\). SP5 filtra la referencia de
movimiento y vuelve a comprobar la acción después del reparto de wrench, la
saturación y la integración. Esta comprobación es necesaria porque una acción
segura antes de limitar actuadores puede dejar de serlo al aplicarse.

\paragraph{Comunicación local.}
La comunicación se representa mediante el grafo
\(\mathcal{G}_c(t)=(\mathcal{R},\mathcal{E}_c(t))\), con
%
\begin{equation}
    \mathcal{E}_c(t)
    =
    \left\{
        \{i,j\}:
        i\neq j,\;
        \|p_i(t)-p_j(t)\|_2\leq R_c
    \right\}.
    \label{eq:communication_graph_methodology}
\end{equation}
%
El robot \(i\) emplea estado propio, percepción disponible y mensajes de
\(\mathcal{N}_i(t)\). Los experimentos modifican el radio \(R_c\), el retardo
entero \(\tau\), la probabilidad de pérdida \(p_{\mathrm{loss}}\), la caducidad de
las copias y la política de retransmisión. Se denomina vecinal a una etapa solo
cuando su estimación, actualización y cierre respetan esa información. Los
registros globales, optimizadores por carga y cierres que consultan toda la
instancia se declaran como dependencias globales.

\paragraph{Rutas y recursos discretos.}
SP7--SP8 utilizan un nivel mesoscópico. Cada coalición selecciona una ruta
\(r_k\) de un catálogo finito y cada ruta ocupa un conjunto
\(\mathcal{E}(r_k)\) de recursos. Para
%
\begin{equation}
    q_e(r)
    =
    \sum_k
    \mathbf{1}\!\left\{e\in\mathcal{E}(r_k)\right\},
    \qquad
    P(r)
    =
    \sum_e
    \binom{q_e(r)}{2},
    \label{eq:route_conflict_methodology}
\end{equation}
%
\(P(r)\) cuenta pares de coaliciones que comparten recursos. Esta abstracción
permite estudiar tráfico, observabilidad y mensajes a mayor escala, pero no
representa interpolación continua, orientación intermedia ni contacto físico.

\subsection{Variables, tratamientos y comparadores}
\label{subsec:variables_comparators}

Los factores independientes cambian por subproblema. Los principales son el
método de coordinación, \(N\), \(K\), la relación entre demanda y flota, la
heterogeneidad de capacidades, la distribución espacial, la geometría de los
contactos, la demanda de wrench, la presencia de obstáculos, el fallo de un
robot, el catálogo de rutas y las condiciones de comunicación. Cada campaña
modifica solo los factores necesarios para responder a su pregunta y mantiene
los restantes bajo una configuración común.

Las variables de respuesta se agrupan en cinco familias. La calidad estratégica
incluye factibilidad, déficit, exceso, completitud, valor y brecha frente a una
referencia. Las variables temporales comprenden formación de coalición,
acoplamiento, entrega, terminación, recuperación y distancia recorrida. Las
variables físicas incluyen error de pose, residual de wrench, saturación, colisión,
holgura y trabajo. La comunicación se mide mediante mensajes, bytes, versiones,
edad de información y conflictos no observados. El coste computacional se
caracteriza mediante CPU, memoria, iteraciones y evaluaciones propias de cada
algoritmo. Las definiciones exactas y sus unidades se fijan en el SP correspondiente.

Los métodos cumplen papeles diferentes. Los optimizadores globales se utilizan
como oráculos, cotas o certificados en instancias acotadas. No representan la
arquitectura operacional propuesta. Los comparadores de referencia pertenecen a
familias compatibles con el modelo del SP. Las heurísticas permiten medir el
coste de una regla más simple. Las ablaciones retiran el cierre, el término de
cuórum, la guardia, la reparación, la reserva o la retransmisión para aislar el
mecanismo asociado con una hipótesis. Los métodos del piloto que adaptan una
idea publicada sin reproducir su algoritmo completo se rotulan como
\emph{proxies} y no heredan sus garantías.

Una solución de optimización solo se denomina óptima cuando la cota se ha
cerrado o la enumeración ha agotado el espacio considerado. Si el tiempo de
resolución termina antes, se informa la mejor solución y su brecha. Del mismo
modo, un método local no se presenta como distribuido si necesita ocupaciones,
ordenamientos, registros o certificados globales.

La relación entre objetivos y evidencia es la siguiente. OE1 se estudia en
SP0--SP2 mediante cierre, equilibrio, completitud y brecha. OE2 se aborda en
SP3--SP4 mediante wrench, contacto, acoplamiento y pose. OE3 se contrasta en
SP5--SP7 y Cargo mediante colisión, bloqueo, reparación y tráfico. OE4 se estudia
en SP8 mediante observabilidad, mensajes, CPU y calidad bajo red imperfecta.
OE5 se evalúa mediante la integración Cargo, la conservación de artefactos y el
piloto Industrial 2.

\subsection{Fases y protocolo experimental}
\label{subsec:experimental_protocol}

El trabajo se desarrolló en cuatro fases. La primera fue la formalización. En cada
SP se definieron el espacio de estados, la información disponible, la acción, las
restricciones, el criterio de éxito y el resultado que contradiría la hipótesis. Las
propiedades de potencial, equilibrio, factibilidad o terminación se demostraron
analíticamente y, cuando el espacio lo permitió, se comprobaron mediante
enumeración de todos los perfiles.

La segunda fase fue de desarrollo. Se implementaron generadores de mundos,
métodos, oráculos, registros e invariantes automáticos. Las semillas de esta fase
se utilizaron para detectar errores, seleccionar rangos de parámetros y comprobar
finitud, restricciones, consistencia dimensional y correspondencia entre
configuraciones y resultados. Estas observaciones no se incorporaron a los
contrastes confirmatorios.

La tercera fase comprendió las campañas por SP. Cada comparación usó mundos
compartidos y parámetros congelados. Los tratamientos partieron del mismo
estado, recibieron los mismos eventos exógenos y conservaron límites, horizonte y
criterio de éxito cuando sus modelos eran compatibles. Las diferencias de
información propias de cada arquitectura se mantuvieron y se documentaron; no se
dio información global a un método vecinal para igualarlo artificialmente con el
oráculo.

La cuarta fase integró la misión Cargo y ejecutó el piloto Industrial 2. Cargo cruzó
cuatro regímenes, tres tamaños de flota y seis tratamientos con treinta semillas
confirmatorias por combinación de régimen y tamaño. Cada mundo incluyó
reclutamiento, acoplamiento, transporte y, según el régimen, obstáculo, red
degradada o fallo. Industrial 2 utilizó cuatro escenarios, cuatro métodos y dos
semillas por escenario durante 60~s. Este tamaño solo permite identificar modos de
fallo y comprobar ejecución funcional.

Solo se excluyen réplicas corruptas, duplicadas o afectadas por un fallo externo al
algoritmo. Las colisiones, excepciones internas, no convergencia, bloqueo y
agotamiento del horizonte permanecen en el denominador. El éxito de misión exige
cierre lógico, factibilidad de la modalidad, ausencia de la colisión definida para la
campaña y entrega dentro del horizonte.

\subsection{Procesamiento estadístico}
\label{subsec:statistical_analysis}

El análisis conserva el emparejamiento. Para dos tratamientos \(A\) y \(B\), la
diferencia del mundo \(w\) es
%
\begin{equation}
    \Delta_w
    =
    Y_A(w)-Y_B(w),
    \qquad
    \widehat{\Delta}
    =
    \frac{1}{n}
    \sum_{w=1}^{n}\Delta_w.
    \label{eq:paired_estimand_methodology}
\end{equation}
%
El signo favorable se fija antes del análisis según la variable. Para tres o más
tratamientos, el mundo actúa como bloque. El remuestreo se realiza sobre mundos
completos, no sobre robots, cargas o muestras temporales.

Cada contraste informa el número de mundos, los fallos, la estimación del efecto,
un intervalo de confianza del 95\,\% y, cuando procede, el valor de \(p\). Los
intervalos \emph{bootstrap} emplean 2000 remuestras pareadas. Las variables
binarias, como éxito, colisión o restauración, se analizan mediante diferencias
pareadas de proporciones y la prueba exacta de McNemar. Las variables continuas
u ordinales se comparan mediante Wilcoxon pareado. Cuando intervienen varios
métodos sobre los mismos mundos se emplean Friedman y un tamaño de efecto
basado en concordancia.

El procedimiento de Holm corrige las familias de contrastes preespecificadas. El
valor de \(p\) no reemplaza la magnitud ni el intervalo del efecto. Un intervalo que
contiene cero se interpreta como ausencia de apoyo concluyente a la diferencia,
no como demostración de equivalencia. Una afirmación de no inferioridad habría
requerido definir previamente un margen clínico o ingenieril, lo que no se hizo en
estas campañas.

Los tiempos de ejecuciones que agotan el horizonte se conservan como valores
truncados o se analizan mediante la regla definida por la campaña. No se sustituyen
por tiempos de entrega. Las medias temporales incondicionales se acompañan de
éxito y causa de terminación para evitar que una colisión temprana parezca una
mejora de velocidad.

El número de semillas se fijó por presupuesto computacional y no mediante un
cálculo de potencia previo. Por ello, la resolución estadística se juzga mediante la
anchura de los intervalos observados. Las conclusiones se restringen a la población
de mundos inducida por los generadores, factores y rangos ensayados.

\subsection{Instrumentos y trazabilidad}
\label{subsec:tools_reproducibility}

Los instrumentos de investigación comprenden los modelos matemáticos, los
simuladores, los solucionadores de optimización, los generadores de mundos, los
registros por ejecución, las pruebas automáticas y los scripts estadísticos. Las
campañas principales se implementaron en Python. Los problemas lineales,
mixtos y cuadráticos se resolvieron con solucionadores numéricos adecuados a
cada formulación. CoppeliaSim se utilizó para el piloto geométrico y cinemático
Industrial 2.

Cada registro conserva, cuando corresponde, identificador del mundo, semilla,
factores, método, parámetros, miembros de las coaliciones, estados de cierre,
certificados, tiempos, errores, eventos de seguridad, mensajes, CPU y causa de
terminación. Las tablas y figuras se generan a partir de esos registros y no de
valores transcritos manualmente.

La reproducibilidad se declara por campaña. SP0--SP1, SP5--SP8, Cargo e
Industrial 2 conservan generadores, configuraciones, semillas y resultados
suficientes para regeneración. SP2--SP4 conservan observaciones, scripts de
postproceso, tablas y auditorías, pero no todos sus generadores históricos. En
estos casos se habla de reanálisis reproducible y no de regeneración completa.
El anexo de reproducibilidad identifica versiones, manifiestos, hashes y
disponibilidad de los artefactos.

\subsection{Validez y límites de inferencia}
\label{subsec:validity}

La validez interna se apoya en mundos pareados, separación entre desarrollo y
confirmación, parámetros congelados y conservación de los fallos en el
denominador. Las ablaciones permiten atribuir efectos a componentes concretos,
aunque una ablación conjunta, como retirar simultáneamente certificado, ruta y
filtro, solo identifica el efecto del bloque completo.

La validez de constructo depende de no confundir variables de capas distintas. La
cardinalidad mide pertenencia, el servicio operacional mide una contribución
normalizada, el residual de wrench mide capacidad mecánica planar, la holgura
mide separación geométrica y el conflicto de rutas mide exclusión lógica. Ninguna
de estas variables sustituye a las restantes. Un residual KKT pequeño tampoco
certifica una asignación entera ni una trayectoria físicamente segura.

La validez de conclusión se protege mediante el emparejamiento, los intervalos de
confianza, la corrección de multiplicidad y la declaración de resultados negativos.
Los contrastes exploratorios y confirmatorios se distinguen. Industrial 2 no se
utiliza para establecer superioridad estadística debido a sus dos semillas por
escenario.

La validez externa queda limitada a los tamaños, distribuciones, geometrías,
contactos y regímenes de comunicación evaluados. El modelo principal es planar,
usa pose conocida, contactos fijos y esfuerzos acotados. No incluye soporte
tridimensional, fricción completa, deformación, percepción real, desgaste,
tracción rueda--suelo ni hardware. Cargo considera una carga activa, un fallo
simple, un líder temporal y dependencias globales declaradas. SP7--SP8 emplean
rutas discretas y no una planta continua.

Por tanto, la integración Cargo demuestra que las capas pueden coexistir dentro
de la planta reducida y los regímenes ensayados. No demuestra estabilidad híbrida
global, seguridad universal, descentralización integral ni rendimiento industrial.
Cada conclusión se mantiene en el estrato donde fue obtenida y solo se transfiere
cuando se conservan su modelo, información, restricciones y criterio de
validación.